\documentclass[12pt, a4paper]{ctexart}
\usepackage{fontspec}
\usepackage{xcolor}
\usepackage{hyperref}
\usepackage{geometry}
\usepackage{enumitem}
\usepackage{parskip}
\usepackage{titlesec}
\usepackage{tcolorbox}
\tcbuselibrary{breakable}
\tcbset{autoparskip}

\usepackage{setspace}
\usepackage{listings}
\usepackage{amssymb}

\geometry{left=2.5cm, right=2.5cm, top=2.5cm, bottom=2.5cm}

\setmainfont{Times New Roman}
\setCJKmainfont{SimSun}

\hypersetup{
    colorlinks=true,
    linkcolor=blue!70!black,
    urlcolor=cyan,
    pdftitle={集合-逐题精讲解义},
    pdfauthor={专升本-fifine老师}
}

\definecolor{myblue}{RGB}{30,100,200}
\definecolor{lightblue}{RGB}{230,240,255}
\definecolor{darkblue}{RGB}{0,50,120}

\newtcolorbox{bluebox}[1][]{
    breakable,
    colback=lightblue,
    colframe=myblue,
    coltitle=darkblue,
    fonttitle=\bfseries,
    title={#1},
    boxrule=0.8pt,
    arc=4pt,
    left=6pt,
    right=6pt,
    top=6pt,
    bottom=6pt,
    before skip=8pt,
    after skip=8pt
}

\usepackage{etoolbox}
\BeforeBeginEnvironment{bluebox}{\par\vfill}%
\AfterEndEnvironment{bluebox}{\par\vfill}%

\titleformat{\section}
  {\normalfont\Large\bfseries\color{myblue}}{\thesection}{1em}{}
\titlespacing*{\section}{0pt}{3.5ex plus 1ex minus .2ex}{1.5ex plus .2ex}

\title{\textcolor{myblue}{\textbf{集合 - 逐题精讲解义}}}
\author{专升本-fifine老师 教学组}
\date{2026年03月05日}

\lstset{
    language=Java,
    basicstyle=\ttfamily\small,
    keywordstyle=\color{blue},
    commentstyle=\color{green!60!black},
    stringstyle=\color{orange},
    numbers=left,
    numberstyle=\tiny\color{gray},
    frame=single,
    breaklines=true,
    columns=fullflexible,
    showstringspaces=false
}

\newcommand{\code}[1]{\texttt{#1}}

\begin{document}

\maketitle
\thispagestyle{empty}
\newpage

\section{集合}

\begin{enumerate}[label=\textbf{\arabic*.}]

	\item \textbf{以下关于Java集合框架的说法正确的是}
	      \begin{enumerate}[label=\Alph*.]
		      \item List允许重复元素
		      \item Set允许重复元素
		      \item Map存储键值对，键可以重复
		      \item 以上都不对
	      \end{enumerate}

	      \begin{bluebox}{答案与深度解析}
		      \textbf{答案：A. List允许重复元素}

		      \textbf{【核心认知框架}
		      \begin{itemize}[label={}, leftmargin=*, nosep]
			      \item \textbf{题目本质}：考查三大集合核心特性对比
			      \item \textbf{关键区分点}：重复元素规则、存储结构、键值唯一性
			      \item \textbf{命题人意图}：测试对集合框架本质差异的理解
		      \end{itemize}

		      \textbf{【选项原理深度拆解】}

		      \begin{itemize}[leftmargin=*, nosep, label={}]
			      \item[A] \textbf{List允许重复元素 — 正确}
			            \begin{itemize}[label={}, leftmargin=*, nosep]
				            \item \textbf{JVM层面实现}：
				                  \begin{lstlisting}[language=Java]
// ArrayList底层使用Object[]数组
transient Object[] elementData;

// add方法直接添加，不检查重复元素
public boolean add(E e) {
    ensureCapacityInternal(size + 1);
    elementData[size++] = e;
    return true;
}
				      \end{lstlisting}
				            \item \textbf{内存结构}：顺序存储，索引访问O(1)
				            \item \textbf{代码验证}：
				                  \begin{lstlisting}[language=Java]
List<String> list = new ArrayList<>();
list.add("A");
list.add("A"); // 允许重复
System.out.println(list); // [A, A]
				      \end{lstlisting}
			            \end{itemize}

			      \item[B] \textbf{Set允许重复元素 — 错误}
			            \begin{itemize}[label={}, leftmargin=*, nosep]
				            \item \textbf{JVM层面实现}：
				                  \begin{lstlisting}[language=Java]
// HashSet依赖HashMap的key唯一性
private transient HashMap<E,Object> map;

public boolean add(E e) {
    return map.put(e, PRESENT) == null;
    // 返回null说明之前不存在，添加成功
    // 返回非null说明之前已存在，添加失败
}
				      \end{lstlisting}
				            \item \textbf{哈希机制}：通过equals()和hashCode()保证唯一性
				            \item \textbf{命题陷阱}：命题人故意混淆List和Set的重复规则
			            \end{itemize}

			      \item[C] \textbf{Map存储键值对，键可以重复 — 错误}
			            \begin{itemize}[label={}, leftmargin=*, nosep]
				            \item \textbf{JVM层面实现}：
				                  \begin{lstlisting}[language=Java]
// HashMap的put方法会覆盖旧值
public V put(K key, V value) {
    return putVal(hash(key), key, value, false, true);
}

final V putVal(int hash, K key, V value,
               boolean onlyIfAbsent, boolean evict) {
    // 如果key已存在，覆盖value
    if (e != null) {
        V oldValue = e.value;
        if (!onlyIfAbsent || oldValue == null)
            e.value = value;
        return oldValue;
    }
}
				      \end{lstlisting}
				            \item \textbf{唯一性保证}：通过哈希冲突处理机制保证key唯一
				            \item \textbf{代码验证}：
				                  \begin{lstlisting}[language=Java]
Map<String, Integer> map = new HashMap<>();
map.put("A", 1);
map.put("A", 2); // 覆盖旧值
System.out.println(map); // {A=2}
				      \end{lstlisting}
			            \end{itemize}

			      \item[D] \textbf{以上都不对 — 错误}
			            \begin{itemize}[label={}, leftmargin=*, nosep]
				            \item A选项陈述正确
				            \item 命题人设置此选项是为了干扰判断
			            \end{itemize}
		      \end{itemize}

		      \textbf{【应背诵知识点}
		      \begin{itemize}[label={}, nosep]
			      \item List：有序、允许重复、可通过索引访问
			      \item Set：无序、不允许重复、依赖equals()和hashCode()
			      \item Map：键值对存储、键唯一、值可重复
			      \item 选项辨析：List=有序重复，Set=无序唯一，Map=键值对
		      \end{itemize}

		      \textbf{【命题趋势与实战策略】}

		      \textbf{考场秒杀}：看到"允许重复"→必选List，"不允许重复"→必选Set
	      \end{bluebox}

	      \vspace{1em}

	\item \textbf{在Java中，以下哪个集合实现了链表结构？}
	      \begin{enumerate}[label=\Alph*.]
		      \item ArrayList
		      \item LinkedList
		      \item HashSet
		      \item HashMap
	      \end{enumerate}

	      \begin{bluebox}{答案与深度解析}
		      \textbf{答案：B. LinkedList}

		      \textbf{【核心认知框架}
		      \begin{itemize}[label={}, leftmargin=*, nosep]
			      \item \textbf{题目本质}：考查集合底层数据结构
			      \item \textbf{关键区分点}：数组结构 vs 链表结构 vs 哈希表结构
			      \item \textbf{命题人意图}：测试对集合实现细节的理解
		      \end{itemize}

		      \textbf{【选项原理深度拆解】}

		      \begin{itemize}[leftmargin=*, nosep, label={}]
			      \item[A] \textbf{ArrayList — 错误（动态数组实现）}
			            \begin{itemize}[label={}, leftmargin=*, nosep]
				            \item \textbf{字节码证据}：
				                  \begin{lstlisting}[language=Java]
// ArrayList声明
public class ArrayList<E> extends AbstractList<E>
    implements List<E>, RandomAccess, Cloneable, java.io.Serializable {

    transient Object[] elementData; // 数组结构
}
				      \end{lstlisting}
				            \item \textbf{内存布局}：连续内存空间，索引访问O(1)
				            \item \textbf{扩容机制}：capacity grows by 1.5倍
				            \item \textbf{适用场景}：频繁随机访问
			            \end{itemize}

			      \item[B] \textbf{LinkedList — 正确（双向链表实现）}
			            \begin{itemize}[label={}, leftmargin=*, nosep]
				            \item \textbf{字节码证据}：
				                  \begin{lstlisting}[language=Java]
// LinkedList声明
public class LinkedList<E>
    extends AbstractSequentialList<E>
    implements List<E>, Deque<E>, Cloneable, java.io.Serializable {

    transient Node<E> first; // 头节点
    transient Node<E> last;  // 尾节点

    // 链表节点定义
    private static class Node<E> {
        E item;
        Node<E> next;
        Node<E> prev;

        Node(Node<E> prev, E element, Node<E> next) {
            this.item = element;
            this.next = next;
            this.prev = prev;
        }
    }
}
				      \end{lstlisting}
				            \item \textbf{双向链表特性}：每个节点有prev和next指针
				            \item \textbf{性能特点}：插入删除O(1)，随机访问O(n)
				            \item \textbf{代码验证}：
				                  \begin{lstlisting}[language=Java]
LinkedList<String> list = new LinkedList<>();
list.add("A"); // 添加到链表尾部
list.addFirst("B"); // 添加到头部
list.addLast("C"); // 添加到尾部
System.out.println(list); // [B, A, C]
				      \end{lstlisting}
			            \end{itemize}

			      \item[C] \textbf{HashSet — 错误（哈希表实现）}
			            \begin{itemize}[label={}, leftmargin=*, nosep]
				            \item \textbf{字节码证据}：
				                  \begin{lstlisting}[language=Java]
public class HashSet<E>
    extends AbstractSet<E>
    implements Set<E>, Cloneable, java.io.Serializable {

    private transient HashMap<E,Object> map;
    // 底层依赖HashMap
}
				      \end{lstlisting}
				            \item \textbf{哈希表结构}：数组+链表+红黑树
				            \item \textbf{查找性能}：O(1)平均
			            \end{itemize}

			      \item[D] \textbf{HashMap — 错误（哈希表实现）}
			            \begin{itemize}[label={}, leftmargin=*, nosep]
				            \item \textbf{字节码证据}：
				                  \begin{lstlisting}[language=Java]
public class HashMap<K,V> extends AbstractMap<K,V>
    implements Map<K,V>, Cloneable, Serializable {

    transient Node<K,V>[] table; // 哈希表
    static final int TREEIFY_THRESHOLD = 8;

    // 哈希桶结构
    static class Node<K,V> implements Map.Entry<K,V> {
        final int hash;
        final K key;
        V value;
        Node<K,V> next;
    }
}
				      \end{lstlisting}
				            \item \textbf{哈希冲突处理}：链表转红黑树（JDK8+）
			            \end{itemize}
		      \end{itemize}

		      \textbf{【应背诵知识点}
		      \begin{itemize}[label={}, nosep]
			      \item ArrayList：动态数组，查询快O(1)，增删慢O(n)
			      \item LinkedList：双向链表，增删快O(1)，查询慢O(n)
			      \item HashSet/HashMap：哈希表，增删改查O(1)
			      \item 选项辨析：LinkedList=链表，ArrayList=数组，HashMap=哈希表
		      \end{itemize}

		      \textbf{【命题趋势与实战策略】}

		      \textbf{考场秒杀}：看到"链表"→选LinkedList
	      \end{bluebox}

	      \vspace{1em}

	\item \textbf{以下哪个集合是无序且不允许重复元素的？}
	      \begin{enumerate}[label=\Alph*.]
		      \item ArrayList
		      \item HashSet
		      \item LinkedHashSet
		      \item TreeSet
	      \end{enumerate}

	      \begin{bluebox}{答案与深度解析}
		      \textbf{答案：B. HashSet}

		      \textbf{【核心认知框架}
		      \begin{itemize}[label={}, leftmargin=*, nosep]
			      \item \textbf{题目本质}：考查集合的有序性和重复性特性
			      \item \textbf{关键区分点}：有序vs无序，允许重复vs不允许重复
			      \item \textbf{命题人意图}：测试对Set子类特性的精确理解
		      \end{itemize}

		      \textbf{【选项原理深度拆解】}

		      \begin{itemize}[leftmargin=*, nosep, label={}]
			      \item[A] \textbf{ArrayList — 错误（有序且允许重复）}
			            \begin{itemize}[label={}, leftmargin=*, nosep]
				            \item \textbf{有序性}：元素按插入顺序存储
				            \item \textbf{重复性}：允许重复元素
				            \item \textbf{代码验证}：
				                  \begin{lstlisting}[language=Java]
ArrayList<String> list = new ArrayList<>();
list.add("A");
list.add("B");
list.add("A");
System.out.println(list); // [A, B, A] 有序允许重复
				      \end{lstlisting}
			            \end{itemize}

			      \item[B] \textbf{HashSet — 正确（无序且不允许重复）}
			            \begin{itemize}[label={}, leftmargin=*, nosep]
				            \item \textbf{字节码证据}：
				                  \begin{lstlisting}[language=Java]
public class HashSet<E> extends AbstractSet<E> {
    private transient HashMap<E,Object> map;
    private static final Object PRESENT = new Object();

    public HashSet() {
        map = new HashMap<>();
    }

    public boolean add(E e) {
        return map.put(e, PRESENT) == null;
    }
}
				      \end{lstlisting}
				            \item \textbf{无序性}：哈希表存储，元素顺序取决于哈希值
				            \item \textbf{不允许重复}：依赖HashMap的key唯一性
				            \item \textbf{代码验证}：
				                  \begin{lstlisting}[language=Java]
HashSet<String> set = new HashSet<>();
set.add("A");
set.add("B");
set.add("A");
System.out.println(set); // [A, B] 或 [B, A] 无序不允许重复
				      \end{lstlisting}
				            \item \textbf{命题陷阱}：虽然顺序不可预测，但"无序"是准确描述
			            \end{itemize}

			      \item[C] \textbf{LinkedHashSet — 错误（有序且不允许重复）}
			            \begin{itemize}[label={}, leftmargin=*, nosep]
				            \item \textbf{字节码证据}：
				                  \begin{lstlisting}[language=Java]
public class LinkedHashSet<E>
    extends HashSet<E>
    implements Set<E>, Cloneable, java.io.Serializable {

    public LinkedHashSet(int initialCapacity, float loadFactor) {
        super(initialCapacity, loadFactor, true);
        // 调用HashSet的accessOrder=true的构造器
    }
}
				      \end{lstlisting}
				            \item \textbf{有序性}：维护插入顺序（LinkedList + HashMap）
				            \item \textbf{不允许重复}：继承自HashSet
				            \item \textbf{代码验证}：
				                  \begin{lstlisting}[language=Java]
LinkedHashSet<String> set = new LinkedHashSet<>();
set.add("A");
set.add("B");
set.add("A");
System.out.println(set); // [A, B] 有序不允许重复
				      \end{lstlisting}
			            \end{itemize}

			      \item[D] \textbf{TreeSet — 错误（有序且不允许重复）}
			            \begin{itemize}[label={}, leftmargin=*, nosep]
				            \item \textbf{字节码证据}：
				                  \begin{lstlisting}[language=Java]
public class TreeSet<E> extends AbstractSet<E>
    implements NavigableSet<E>, Cloneable, java.io.Serializable {

    private transient NavigableMap<E,Object> m;

    public TreeSet() {
        this(new TreeMap<>());
    }
}
				      \end{lstlisting}
				            \item \textbf{有序性}：自然排序或自定义排序（红黑树）
				            \item \textbf{不允许重复}：红黑树的BST特性
				            \item \textbf{代码验证}：
				                  \begin{lstlisting}[language=Java]
TreeSet<String> set = new TreeSet<>();
set.add("C");
set.add("A");
set.add("B");
System.out.println(set); // [A, B, C] 排序不允许重复
				      \end{lstlisting}
			            \end{itemize}
		      \end{itemize}

		      \textbf{【应背诵知识点}
		      \begin{itemize}[label={}, nosep]
			      \item HashSet：无序、不允许重复、基于哈希表
			      \item LinkedHashSet：有序（插入顺序）、不允许重复、哈希表+链表
			      \item TreeSet：有序（排序）、不允许重复、红黑树
			      \item 选项辨析：HashSet=无序唯一，LinkedHashSet=有序唯一，TreeSet=排序唯一
		      \end{itemize}

		      \textbf{【命题趋势与实战策略】}

		      \textbf{考场秒杀}：看到"无序且不允许重复"→选HashSet
	      \end{bluebox}

	      \vspace{1em}

	\item \textbf{以下关于HashMap的说法错误的是}
	      \begin{enumerate}[label=\Alph*.]
		      \item 存储键值对
		      \item 键不允许重复
		      \item 是线程安全的
		      \item 基于哈希表实现
	      \end{enumerate}

	      \begin{bluebox}{答案与深度解析}
		      \textbf{答案：C. 是线程安全的}

		      \textbf{【核心认知框架}
		      \begin{itemize}[label={}, leftmargin=*, nosep]
			      \item \textbf{题目本质}：考查HashMap的线程安全性
			      \item \textbf{关键区分点}：线程安全 vs 线程不安全
			      \item \textbf{命题人意图}：测试对并发编程基础的理解
		      \end{itemize}

		      \textbf{【选项原理深度拆解】}

		      \begin{itemize}[leftmargin=*, nosep, label={}]
			      \item[A] \textbf{存储键值对 — 正确}
			            \begin{itemize}[label={}, leftmargin=*, nosep]
				            \item \textbf{字节码证据}：
				                  \begin{lstlisting}[language=Java]
public class HashMap<K,V> extends AbstractMap<K,V> {
    public V put(K key, V value) {
        return putVal(hash(key), key, value, false, true);
    }

    public V get(Object key) {
        Node<K,V> e;
        return (e = getNode(hash(key), key)) == null ? null : e.value;
    }
}
				      \end{lstlisting}
				            \item \textbf{存储结构}：键值对 Entry<K,V>
				            \item \textbf{代码验证}：
				                  \begin{lstlisting}[language=Java]
Map<String, Integer> map = new HashMap<>();
map.put("name", 18);
System.out.println(map.get("name")); // 18
				      \end{lstlisting}
			            \end{itemize}

			      \item[B] \textbf{键不允许重复 — 正确}
			            \begin{itemize}[label={}, leftmargin=*, nosep]
				            \item \textbf{哈希冲突处理}：
				                  \begin{lstlisting}[language=Java]
final V putVal(int hash, K key, V value, ...) {
    Node<K,V>[] tab; Node<K,V> p; int n, i;
    // 通过哈希值确定索引
    if ((p = tab[i = (n - 1) & hash]) == null)
        tab[i] = newNode(hash, key, value, null);
    else {
        // key已存在，覆盖value
        if (e.hash == hash &&
            ((k = e.key) == key || (key != null && key.equals(k)))) {
            e.value = value;
        }
    }
}
				      \end{lstlisting}
				            \item \textbf{唯一性保证}：通过equals()和hashCode()
			            \end{itemize}

			      \item[C] \textbf{是线程安全的 — 错误（本题答案）}
			            \begin{itemize}[label={}, leftmargin=*, nosep]
				            \item \textbf{线程不安全证据}：
				                  \begin{lstlisting}[language=Java]
public V put(K key, V value) {
    // 没有任何同步机制
    return putVal(hash(key), key, value, false, true);
}

// 扩容操作不是原子的
final Node<K,V>[] resize() {
    // 多线程扩容会导致死循环或数据丢失
}
				      \end{lstlisting}
				            \item \textbf{并发问题}：
				                  \begin{itemize}[label={}, nosep]
					                  \item 数据覆盖：多线程put同一位置
					                  \item 扩容死循环：链表环化（JDK7）
					                  \item 丢失数据：resize时覆盖旧数据
				                  \end{itemize}
				            \item \textbf{线程安全替代方案}：
				                  \begin{itemize}[label={}, nosep]
					                  \item Hashtable（过时，全表锁）
					                  \item ConcurrentHashMap（推荐，分段锁）
					                  \item Collections.synchronizedMap()
				                  \end{itemize}
				            \item \textbf{代码验证（并发问题）}：
				                  \begin{lstlisting}[language=Java]
Map<String, Integer> map = new HashMap<>();
Runnable task = () -> {
    for (int i = 0; i < 1000; i++) {
        map.put("key", map.getOrDefault("key", 0) + 1);
    }
};
// 多线程执行后，值可能小于预期
				      \end{lstlisting}
				            \item \textbf{命题陷阱}：这是高频考点，常考HashMap线程不安全
			            \end{itemize}

			      \item[D] \textbf{基于哈希表实现 — 正确}
			            \begin{itemize}[label={}, leftmargin=*, nosep]
				            \item \textbf{数据结构}：
				                  \begin{lstlisting}[language=Java]
transient Node<K,V>[] table; // 哈希表

// JDK8+：数组 + 链表 + 红黑树
static final int TREEIFY_THRESHOLD = 8;
static final int UNTREEIFY_THRESHOLD = 6;
				      \end{lstlisting}
				            \item \textbf{性能特性}：增删改查平均O(1)
			            \end{itemize}
		      \end{itemize}

		      \textbf{【应背诵知识点}
		      \begin{itemize}[label={}, nosep]
			      \item HashMap：线程不安全、基于哈希表、键唯一
			      \item ConcurrentHashMap：线程安全、分段锁
			      \item Hashtable：线程安全、全表锁（性能差）
			      \item 选项辨析：HashMap=不安全，ConcurrentHashMap=安全
		      \end{itemize}

		      \textbf{【命题趋势与实战策略】}

		      \textbf{考场秒杀}：看到HashMap选"线程安全"→必错，选"线程不安全"→必对
	      \end{bluebox}

	      \vspace{1em}

	\item \textbf{以下哪个方法可以判断一个集合是否为空？}
	      \begin{enumerate}[label=\Alph*.]
		      \item isEmpty()
		      \item isBlank()
		      \item isNullOrEmpty()
		      \item isNotNull()
	      \end{enumerate}

	      \begin{bluebox}{答案与深度解析}
		      \textbf{答案：A. isEmpty()}

		      \textbf{【核心认知框架}
		      \begin{itemize}[label={}, leftmargin=*, nosep]
			      \item \textbf{题目本质}：考查Collection接口的标准方法
			      \item \textbf{关键区分点}：标准接口方法 vs 第三方库方法
			      \item \textbf{命题人意图}：测试对Java集合API的掌握
		      \end{itemize}

		      \textbf{【选项原理深度拆解】}

		      \begin{itemize}[leftmargin=*, nosep, label={}]
			      \item[A] \textbf{isEmpty() — 正确}
			            \begin{itemize}[label={}, leftmargin=*, nosep]
				            \item \textbf{字节码证据}：
				                  \begin{lstlisting}[language=Java]
public interface Collection<E> extends Iterable<E> {
    // Collection接口中定义的标准方法
    boolean isEmpty();
}

// ArrayList实现
public boolean isEmpty() {
    return size == 0;
}

// HashSet实现
public boolean isEmpty() {
    return map.isEmpty();
}
				      \end{lstlisting}
				            \item \textbf{标准API}：所有集合通用方法
				            \item \textbf{代码验证}：
				                  \begin{lstlisting}[language=Java]
List<String> list = new ArrayList<>();
System.out.println(list.isEmpty()); // true

list.add("A");
System.out.println(list.isEmpty()); // false
				      \end{lstlisting}
			            \end{itemize}

			      \item[B] \textbf{isBlank() — 错误（非集合方法）}
			            \begin{itemize}[label={}, leftmargin=*, nosep]
				            \item \textbf{实际来源}：Apache Commons Lang或Spring工具类
				            \item \textbf{常见用法}：
				                  \begin{lstlisting}[language=Java]
// Apache Commons Lang
org.apache.commons.lang3.StringUtils.isBlank(String str)

// Spring
org.springframework.util.StringUtils.hasText(String str)
				      \end{lstlisting}
				            \item \textbf{命题陷阱}：不是Java标准集合API
			            \end{itemize}

			      \item[C] \textbf{isNullOrEmpty() — 错误（非集合方法）}
			            \begin{itemize}[label={}, leftmargin=*, nosep]
				            \item \textbf{实际来源}：Apache Commons Lang
				                  \begin{lstlisting}[language=Java]
// Apache Commons Lang
org.apache.commons.lang3.CollectionUtils.isEmpty(Collection<?> coll)
org.apache.commons.lang3.StringUtils.isEmpty(String str)
				      \end{lstlisting}
				            \item \textbf{注意事项}：不是JDK标准API
			            \end{itemize}

			      \item[D] \textbf{isNotNull() — 错误（非集合方法）}
			            \begin{itemize}[label={}, leftmargin=*, nosep]
				            \item \textbf{实际来源}：断言库或测试框架
				                  \begin{lstlisting}[language=Java]
// Apache Commons Lang
org.apache.commons.lang3.Validate.notNull(Object reference)

// Guava
com.google.common.base.Preconditions.checkNotNull(Object reference)
				      \end{lstlisting}
				            \item \textbf{正确判断}：
				                  \begin{lstlisting}[language=Java]
// 判断集合不为空
if (list != null && !list.isEmpty()) {
    // 处理集合
}
				      \end{lstlisting}
			            \end{itemize}
		      \end{itemize}

		      \textbf{【应背诵知识点}
		      \begin{itemize}[label={}, nosep]
			      \item isEmpty()：Collection接口标准方法
			      \item size()：返回元素个数
			      \item 判断非空：list != null && !list.isEmpty()
			      \item 选项辨析：isEmpty()=JDK标准，isBlank/isNullOrEmpty=第三方库
		      \end{itemize}

		      \textbf{【命题趋势与实战策略】}

		      \textbf{考场秒杀}：看到"判断集合是否为空"→选isEmpty()
	      \end{bluebox}

	      \vspace{1em}

	\item \textbf{以下关于Java中Lambda表达式的说法正确的是}
	      \begin{enumerate}[label=\Alph*.]
		      \item 可以简化匿名内部类的写法
		      \item 不支持函数式接口
		      \item 不能作为方法参数
		      \item 性能比普通方法差
	      \end{enumerate}

	      \begin{bluebox}{答案与深度解析}
		      \textbf{答案：A. 可以简化匿名内部类的写法}

		      \textbf{【核心认知框架}
		      \begin{itemize}[label={}, leftmargin=*, nosep]
			      \item \textbf{题目本质}：考查Lambda表达式的核心价值
			      \item \textbf{关键区分点}：语法糖本质 vs 功能限制
			      \item \textbf{命题人意图}：测试对Lambda的设计目的理解
		      \end{itemize}

		      \textbf{【选项原理深度拆解】}

		      \begin{itemize}[leftmargin=*, nosep, label={}]
			      \item[A] \textbf{可以简化匿名内部类的写法 — 正确}
			            \begin{itemize}[label={}, leftmargin=*, nosep]
				            \item \textbf{代码对比}：
				                  \begin{lstlisting}[language=Java]
// 传统匿名内部类
Runnable r1 = new Runnable() {
    @Override
    public void run() {
        System.out.println("Hello");
    }
};

// Lambda表达式
Runnable r2 = () -> System.out.println("Hello");

// List遍历对比
// 传统方式
for (String s : list) {
    System.out.println(s);
}

// Lambda方式
list.forEach(s -> System.out.println(s));
				      \end{lstlisting}
				            \item \textbf{字节码证据}：
				                  \begin{lstlisting}[language=Java]
// Lambda编译为私有静态方法或私有实例方法
// 运行时通过invokedynamic指令动态调用
// 比匿名内部类生成的类文件更少
				      \end{lstlisting}
				            \item \textbf{简化程度}：代码量减少50~80%
			            \end{itemize}

			      \item[B] \textbf{不支持函数式接口 — 错误}
			            \begin{itemize}[label={}, leftmargin=*, nosep]
				            \item \textbf{反证}：Lambda只能用于函数式接口
				                  \begin{lstlisting}[language=Java]
@FunctionalInterface
interface Calculator {
    int calculate(int a, int b);
}

// Lambda表达式
Calculator add = (a, b) -> a + b;
Calculator multi = (a, b) -> a * b;

System.out.println(add.calculate(3, 5)); // 8
				      \end{lstlisting}
				            \item \textbf{函数式接口定义}：只包含一个抽象方法的接口
				            \item \textbf{常见函数式接口}：
				                  \begin{itemize}[label={}, nosep]
					                  \item Runnable：无参无返回
					                  \item Callable：无参有返回
					                  \item Consumer：有参无返回
					                  \item Function：有参有返回
					                  \item Predicate：有参返回boolean
				                  \end{itemize}
			            \end{itemize}

			      \item[C] \textbf{不能作为方法参数 — 错误}
			            \begin{itemize}[label={}, leftmargin=*, nosep]
				            \item \textbf{代码反证}：
				                  \begin{lstlisting}[language=Java]
// Lambda作为方法参数
public void process(List<Integer> list,
                    Consumer<Integer> consumer) {
    list.forEach(consumer);
}

// 调用
List<Integer> numbers = Arrays.asList(1, 2, 3);
process(numbers, n -> System.out.println(n * 2));
				      \end{lstlisting}
				            \item \textbf{实际应用}：
				                  \begin{itemize}[label={}, nosep]
					                  \item Stream API：filter、map、forEach
					                  \item 线程池：submit(() -> {})
					                  \item 事件处理：onClick(e -> {})
				                  \end{itemize}
			            \end{itemize}

			      \item[D] \textbf{性能比普通方法差 — 错误}
			            \begin{itemize}[label={}, leftmargin=*, nosep]
				            \item \textbf{性能对比}：
				                  \begin{itemize}[label={}, nosep]
					                  \item Lambda性能接近普通方法
					                  \item 匿名内部类需要生成类文件，可能有额外开销
					                  \item invokedynamic指令优化了调用性能
				                  \end{itemize}
				            \item \textbf{JDK优化}：
				                  \begin{lstlisting}[language=Java]
// JDK8+: Lambda被编译为调用点处的invokedynamic指令
// 运行时动态绑定为方法句柄
// 性能与直接调用方法相近
				      \end{lstlisting}
				            \item \textbf{实测数据}：Lambda性能 ≤ 普通方法 < 匿名内部类
			            \end{itemize}
		      \end{itemize}

		      \textbf{【应背诵知识点}
		      \begin{itemize}[label={}, nosep]
			      \item Lambda：简化匿名内部类、用于函数式接口、可作方法参数
			      \item 函数式接口：只有一个抽象方法的接口
			      \item 常见接口：Runnable、Consumer、Function、Predicate
			      \item 语法：(参数) -> {方法体} 或 参数 -> {方法体} 或 参数 -> 表达式
		      \end{itemize}

		      \textbf{【命题趋势与实战策略】}

		      \textbf{考场秒杀}：看到Lambda"简化"→选正确，"不支持/不能"→选错误
	      \end{bluebox}

	      \vspace{1em}

	\item \textbf{以下哪个是函数式接口？}
	      \begin{enumerate}[label=\Alph*.]
		      \item 接口中有多个抽象方法
		      \item 接口中有一个抽象方法
		      \item 接口中有静态方法
		      \item 接口中有默认方法
	      \end{enumerate}

	      \begin{bluebox}{答案与深度解析}
		      \textbf{答案：B. 接口中有一个抽象方法}

		      \textbf{【核心认知框架}
		      \begin{itemize}[label={}, leftmargin=*, nosep]
			      \item \textbf{题目本质}：考查函数式接口的精确定义
			      \item \textbf{关键区分点}：抽象方法数量是唯一标准
			      \item \textbf{命题人意图}：测试对函数式接口规范的理解
		      \end{itemize}

		      \textbf{【选项原理深度拆解】}

		      \begin{itemize}[leftmargin=*, nosep, label={}]
			      \item[A] \textbf{接口中有多个抽象方法 — 错误}
			            \begin{itemize}[label={}, leftmargin=*, nosep]
				            \item \textbf{反证}：多抽象方法不能用于Lambda
				                  \begin{lstlisting}[language=Java]
// 非函数式接口（多个抽象方法）
interface NotFunctional {
    void method1();
    void method2();
}

// 编译错误：目标类型不是函数式接口
// Runnable r = () -> {}; // OK
// NotFunctional nf = () -> {}; // 错误
				      \end{lstlisting}
				            \item \textbf{Lambda限制}：只能用于函数式接口
			            \end{itemize}

			      \item[B] \textbf{接口中有一个抽象方法 — 正确}
			            \begin{itemize}[label={}, leftmargin=*, nosep]
				            \item \textbf{定义依据}：JLS §9.8
				            \item \textbf{@FunctionalInterface注解}：
				                  \begin{lstlisting}[language=Java]
@FunctionalInterface
interface Calculator {
    int calculate(int a, int b);
}

// 编译器会检查是否只有一个抽象方法
@FunctionalInterface
interface Functional {
    void abstractMethod();
    static void staticMethod() {}    // OK
    default void defaultMethod() {}  // OK
}

// 错误：多个抽象方法
@FunctionalInterface  // 编译错误
interface Invalid {
    void method1();
    void method2();
}
				      \end{lstlisting}
				            \item \textbf{重要说明}：
				                  \begin{itemize}[label={}, nosep]
					                  \item 可以有多个默认方法
					                  \item 可以有多个静态方法
					                  \item 可以继承Object的方法
					                  \item 但只能有一个自己的抽象方法
				                  \end{itemize}
			            \end{itemize}

			      \item[C] \textbf{接口中有静态方法 — 错误}
			            \begin{itemize}[label={}, leftmargin=*, nosep]
				            \item \textbf{不影响判断}：静态方法不是抽象方法
				                  \begin{lstlisting}[language=Java]
@FunctionalInterface
interface FunctionalWithStatic {
    void abstractMethod();

    static String getName() {
        return "functional";
    }
}

// OK：仍然是函数式接口
FunctionalWithStatic fs = () -> {};
				      \end{lstlisting}
				            \item \textbf{命题陷阱}：静?方法不影响函数式接口判断
			            \end{itemize}

			      \item[D] \textbf{接口中有默认方法 — 错误}
			            \begin{itemize}[label={}, leftmargin=*, nosep]
				            \item \textbf{不影响判断}：默认方法不是抽象方法
				                  \begin{lstlisting}[language=Java]
@FunctionalInterface
interface FunctionalWithDefault {
    void abstractMethod();

    default void greet() {
        System.out.println("Hello");
    }

    default void farewell() {
        System.out.println("Goodbye");
    }
}

// OK：仍然是函数式接口
FunctionalWithDefault fd = () -> {};
				      \end{lstlisting}
				            \item \textbf{Java8特性}：默认方法让接口可扩展
			            \end{itemize}
		      \end{itemize}

		      \textbf{【应背诵知识点}
		      \begin{itemize}[label={}, nosep]
			      \item 函数式接口：只有一个抽象方法的接口
			      \item 可包含：多个默认方法、多个静态方法
			      \item @FunctionalInterface：可选，编译器检查
			      \item 常见接口：Runnable、Consumer、Function、Predicate、Supplier
		      \end{itemize}

		      \textbf{【命题趋势与实战策略】}

		      \textbf{考场秒杀}：看到函数式接口定义→选"一个抽象方法"，静态/默认方法选错误
	      \end{bluebox}

	      \vspace{1em}

	\item \textbf{以下关于Java中方法引用的说法错误的是}
	      \begin{enumerate}[label=\Alph*.]
		      \item 可以减少代码量
		      \item 有四种类型的方法引用
		      \item 不能引用静态方法
		      \item 是Lambda表达式的一种简洁写法
	      \end{enumerate}

	      \begin{bluebox}{答案与深度解析}
		      \textbf{答案：C. 不能引用静态方法}

		      \textbf{【核心认知框架}
		      \begin{itemize}[label={}, leftmargin=*, nosep]
			      \item \textbf{题目本质}：考查方法引用的语法和类型
			      \item \textbf{关键区分点}：四种类型的能力范围
			      \item \textbf{命题人意图}：测试对方法引用完整语法的掌握
		      \end{itemize}

		      \textbf{【选项原理深度拆解】}

		      \begin{itemize}[leftmargin=*, nosep, label={}]
			      \item[A] \textbf{可以减少代码量 — 正确}
			            \begin{itemize}[label={}, leftmargin=*, nosep]
				            \item \textbf{代码对比}：
				                  \begin{lstlisting}[language=Java]
// Lambda表达式
list.forEach(s -> System.out.println(s));

// 方法引用
list.forEach(System.out::println);

// Lambda
list.stream().map(s -> s.toUpperCase());

// 方法引用
list.stream().map(String::toUpperCase);
				      \end{lstlisting}
				            \item \textbf{简化程度}：平均减少30~50%代码量
			            \end{itemize}

			      \item[B] \textbf{有四种类型的方法引用 — 正确}
			            \begin{itemize}[label={}, leftmargin=*, nosep]
				            \item \textbf{四种类型}：
				                  \begin{lstlisting}[language=Java]
// 1. 静态方法引用
Function<String, Integer> func1 = Integer::parseInt;
Integer value1 = func1.apply("123");

// 2. 对象实例方法引用
Consumer<String> consumer = System.out::println;
consumer.accept("Hello");

// 3. 类实例方法引用
BiPredicate<String, String> bp = String::equals;
boolean result = bp.test("abc", "abc");

// 4. 构造器引用
Supplier<List<String>> supplier = ArrayList::new;
List<String> list = supplier.get();
				      \end{lstlisting}
				            \item \textbf{语法总结}：
				                  \begin{itemize}[label={}, nosep]
					                  \item ClassName::staticMethod
					                  \item object::instanceMethod
					                  \item ClassName::instanceMethod
					                  \item ClassName::new
				                  \end{itemize}
			            \end{itemize}

			      \item[C] \textbf{不能引用静态方法 — 错误（本题答案）}
			            \begin{itemize}[label={}, leftmargin=*, nosep]
				            \item \textbf{代码反证}：
				                  \begin{lstlisting}[language=Java]
// 静态方法引用
List<String> list = Arrays.asList("1", "2", "3");

// Lambda方式
list.stream().map(s -> Integer.parseInt(s));

// 方法引用方式
list.stream().map(Integer::parseInt)
           .forEach(System.out::println);

// 输出：1 2 3
				      \end{lstlisting}
				            \item \textbf{更多示例}：
				                  \begin{lstlisting}[language=Java]
// Math类静态方法
UnaryOperator<Double> sqrt = Math::sqrt;
System.out.println(sqrt.apply(16.0)); // 4.0

// Arrays类静态方法
Function<String[], List<String>> toList = Arrays::asList;
List<String> result = toList.apply(new String[]{"A", "B"});
				      \end{lstlisting}
				            \item \textbf{命题陷阱}：这是方法引用最常用的类型之一
			            \end{itemize}

			      \item[D] \textbf{是Lambda表达式的一种简洁写法 — 正确}
			            \begin{itemize}[label={}, leftmargin=*, nosep]
				            \item \textbf{等价关系}：
				                  \begin{lstlisting}[language=Java]
// Lambda
Function<String, Integer> f1 = s -> Integer.parseInt(s);

// 方法引用（等价）
Function<String, Integer> f2 = Integer::parseInt;

// 编译后字节码相同
				      \end{lstlisting}
				            \item \textbf{适用条件}：
				                  \begin{itemize}[label={}, nosep]
					                  \item Lambda体只调用一个方法
					                  \item 方法参数和返回类型匹配函数式接口
				                  \end{itemize}
			            \end{itemize}
		      \end{itemize}

		      \textbf{【应背诵知识点}
		      \begin{itemize}[label={}, nosep]
			      \item 方法引用：Lambda的简洁写法
			      \item 四种类型：静态方法、实例方法、类实例方法、构造器
			      \item 语法：ClassName::method 或 object::method
			      \item 优势：减少代码量、提高可读性
		      \end{itemize}

		      \textbf{【命题趋势与实战策略】}

		      \textbf{考场秒杀}：看到"不能引用静态方法"→必错，方法引用完全支持静态方法
	      \end{bluebox}

	      \vspace{1em}

	\item \textbf{以下关于Java中Stream流的说法正确的是}
	      \begin{enumerate}[label=\Alph*.]
		      \item 可以对集合进行高效的操作
		      \item 只能顺序处理数据
		      \item 不支持并行处理
		      \item 不能进行数据过滤
	      \end{enumerate}

	      \begin{bluebox}{答案与深度解析}
		      \textbf{答案：A. 可以对集合进行高效的操作}

		      \textbf{【核心认知框架}
		      \begin{itemize}[label={}, leftmargin=*, nosep]
			      \item \textbf{题目本质}：考查Stream API的核心价值
			      \item \textbf{关键区分点}：声明式编程 vs 命令式编程
			      \item \textbf{命题人意图}：测试对Stream设计理念的理解
		      \end{itemize}

		      \textbf{【选项原理深度拆解】}

		      \begin{itemize}[leftmargin=*, nosep, label={}]
			      \item[A] \textbf{可以对集合进行高效的操作 — 正确}
			            \begin{itemize}[label={}, leftmargin=*, nosep]
				            \item \textbf{代码对比}：
				                  \begin{lstlisting}[language=Java]
// 传统方式（命令式）
List<String> result = new ArrayList<>();
for (String s : names) {
    if (s.length() > 3) {
        result.add(s.toUpperCase());
    }
}

// Stream方式（声明式）
List<String> result = names.stream()
    .filter(s -> s.length() > 3)
    .map(String::toUpperCase)
    .collect(Collectors.toList());
				      \end{lstlisting}
				            \item \textbf{高效性体现}：
				                  \begin{itemize}[label={}, nosep]
					                  \item 内部迭代（由Stream控制）
					                  \item 惰性求值（中间操作延迟执行）
					                  \item 短路优化（limit等提前终止）
					                  \item 并行优化（parallelStream）
				                  \end{itemize}
				            \item \textbf{常见操作}：filter、map、sorted、limit、distinct
			            \end{itemize}

			      \item[B] \textbf{只能顺序处理数据 — 错误}
			            \begin{itemize}[label={}, leftmargin=*, nosep]
				            \item \textbf{代码反证}：
				                  \begin{lstlisting}[language=Java]
List<Integer> list = Arrays.asList(1, 2, 3, 4, 5);

// 顺序流
long start1 = System.nanoTime();
long count1 = list.stream().filter(n -> n > 2).count();
long end1 = System.nanoTime();

// 并行流
long start2 = System.nanoTime();
long count2 = list.parallelStream().filter(n -> n > 2).count();
long end2 = System.nanoTime();
				      \end{lstlisting}
				            \item \textbf{并行流特性}：
				                  \begin{itemize}[label={}, nosep]
					                  \item 自动拆分任务（ForkJoinPool）
					                  \item 多线程并行执行
					                  \item 大数据集性能提升明显
				                  \end{itemize}
			            \end{itemize}

			      \item[C] \textbf{不支持并行处理 — 错误}
			            \begin{itemize}[label={}, leftmargin=*, nosep]
				            \item \textbf{并行流API}：
				                  \begin{lstlisting}[language=Java]
// 创建并行流的方式
IntStream.range(1, 1000000).parallel()

// 从集合创建
list.parallelStream()

// 配置线程池
System.setProperty("java.util.concurrent.ForkJoinPool.common.parallelism", "4");
				      \end{lstlisting}
				            \item \textbf{性能对比}：
				                  \begin{itemize}[label={}, nosep]
					                  \item 小数据集：并行可能更慢（线程开销）
					                  \item 大数据集：并行可提升2~4倍性能
					                  \item CPU密集型：并行效率高
				                  \end{itemize}
			            \end{itemize}

			      \item[D] \textbf{不能进行数据过滤 — 错误}
			            \begin{itemize}[label={}, leftmargin=*, nosep]
				            \item \textbf{代码反证}：
				                  \begin{lstlisting}[language=Java]
List<Integer> numbers = Arrays.asList(1, 2, 3, 4, 5, 6);

// 过滤偶数
List<Integer> evens = numbers.stream()
    .filter(n -> n % 2 == 0)
    .collect(Collectors.toList());

// 输出：[2, 4, 6]
System.out.println(evens);

// 复杂过滤
List<String> filtered = names.stream()
    .filter(s -> s.startsWith("A"))
    .filter(s -> s.length() > 3)
    .distinct()
    .collect(Collectors.toList());
				      \end{lstlisting}
				            \item \textbf{过滤操作}：filter、distinct、limit、skip
				            \item \textbf{命题陷阱}：filter是Stream最常用的操作之一
			            \end{itemize}
		      \end{itemize}

		      \textbf{【应背诵知识点}
		      \begin{itemize}[label={}, nosep]
			      \item Stream：声明式集合操作、支持顺序和并行
			      \item 中间操作：filter、map、sorted、distinct
			      \item 终端操作：collect、forEach、count、reduce
			      \item 优势：代码简洁、可并行、惰性求值
		      \end{itemize}

		      \textbf{【命题趋势与实战策略】}

		      \textbf{考场秒杀}：看到Stream"不能"→必错，Stream功能很强大
	      \end{bluebox}

	      \vspace{1em}

	\item \textbf{以下哪个方法用于创建Stream流？}
	      \begin{enumerate}[label=\Alph*.]
		      \item stream()
		      \item createStream()
		      \item makeStream()
		      \item generateStream()
	      \end{enumerate}

	      \begin{bluebox}{答案与深度解析}
		      \textbf{答案：A. stream()}

		      \textbf{【核心认知框架}
		      \begin{itemize}[label={}, leftmargin=*, nosep]
			      \item \textbf{题目本质}：考查Stream API的标准创建方法
			      \item \textbf{关键区分点}：标准API vs 虚构方法名
			      \item \textbf{命题人意图}：测试对Java Stream API的掌握
		      \end{itemize}

		      \textbf{【选项原理深度拆解】}

		      \begin{itemize}[leftmargin=*, nosep, label={}]
			      \item[A] \textbf{stream() — 正确}
			            \begin{itemize}[label={}, leftmargin=*, nosep]
				            \item \textbf{字节码证据}：
				                  \begin{lstlisting}[language=Java]
// Collection接口中定义
public interface Collection<E> extends Iterable<E> {
    default Stream<E> stream() {
        return StreamSupport.stream(spliterator(), false);
    }
}

// List实现
public interface List<E> extends Collection<E> {
    // 继承stream()方法
}

// ArrayList实现（无需重写，使用默认实现）
				      \end{lstlisting}
				            \item \textbf{代码验证}：
				                  \begin{lstlisting}[language=Java]
// 从集合创建
List<String> list = Arrays.asList("A", "B", "C");
Stream<String> stream1 = list.stream();

// 从Set创建
Set<Integer> set = new HashSet<>();
set.add(1);
set.add(2);
Stream<Integer> stream2 = set.stream();

// 并行流
Stream<Integer> parallelStream1 = set.parallelStream();
				      \end{lstlisting}
				            \item \textbf{其他创建方式}：
				                  \begin{lstlisting}[language=Java]
// 从数组创建
Stream<String> stream3 = Arrays.stream(new String[]{"X", "Y", "Z"});

// 从值创建
Stream<String> stream4 = Stream.of("A", "B", "C");

// 范围创建
IntStream range1 = IntStream.range(1, 10); // 1~9
IntStream range2 = IntStream.rangeClosed(1, 10); // 1~10

// 生成创建
Stream<Double> randoms = Stream.generate(Math::random).limit(10);
				      \end{lstlisting}
			            \end{itemize}

			      \item[B] \textbf{createStream() — 错误}
			            \begin{itemize}[label={}, leftmargin=*, nosep]
				            \item \textbf{虚构API}：Java标准库中没有此方法
				            \item \textbf{常见误区}：与Spring或其他框架混淆
				            \item \textbf{命题陷阱}：看似合理但不存在的方法
			            \end{itemize}

			      \item[C] \textbf{makeStream() — 错误}
			            \begin{itemize}[label={}, leftmargin=*, nosep]
				            \item \textbf{虚构API}：Java标准库中没有此方法
				            \item \textbf{正确方法}：应该使用stream()或Stream.of()
			            \end{itemize}

			      \item[D] \textbf{generateStream() — 错误}
			            \begin{itemize}[label={}, leftmargin=*, nosep]
				            \item \textbf{相近API}：Stream.generate()
				                  \begin{lstlisting}[language=Java]
// 正确用法
Stream.generate(() -> "Hello").limit(5).forEach(System.out::println);
				      \end{lstlisting}
				            \item \textbf{方法混淆}：generate是Stream的静态方法，不是实例方法
			            \end{itemize}
		      \end{itemize}

		      \textbf{【应背诵知识点}
		      \begin{itemize}[label={}, nosep]
			      \item Collection.stream()：从创建Stream
			      \item Stream.of()：从值创建Stream
			      \item Arrays.stream()：从数组创建Stream
			      \item Stream.generate()：生成无限Stream
			      \item parallelStream()：创建并行Stream
		      \end{itemize}

		      \textbf{【命题趋势与实战策略】}

		      \textbf{考场秒杀}：看到"创建Stream流"→选stream()或Stream.of()
	      \end{bluebox}

	      \vspace{1em}

	\item \textbf{以下关于Java中Stream流的终端操作的说法错误的是}
	      \begin{enumerate}[label=\Alph*.]
		      \item 终端操作会触发流的执行
		      \item forEach()是终端操作
		      \item 终端操作只能执行一次
		      \item 终端操作后可以继续进行中间操作
	      \end{enumerate}

	      \begin{bluebox}{答案与深度解析}
		      \textbf{答案：D. 终端操作后可以继续进行中间操作}

		      \textbf{【核心认知框架}
		      \begin{itemize}[label={}, leftmargin=*, nosep]
			      \item \textbf{题目本质}：考查Stream的生命周期特性
			      \item \textbf{关键区分点}：中间操作 vs 终端操作
			      \item \textbf{命题人意图}：测试对Stream消费机制的理解
		      \end{itemize}

		      \textbf{【选项原理深度拆解】}

		      \begin{itemize}[leftmargin=*, nosep, label={}]
			      \item[A] \textbf{终端操作会触发流的执行 — 正确}
			            \begin{itemize}[label={}, leftmargin=*, nosep]
				            \item \textbf{字节码证据}：
				                  \begin{lstlisting}[language=Java]
// 中间操作：惰性求值，不立即执行
Stream<String> stream = list.stream()
    .filter(s -> {
        System.out.println("Filtering: " + s);
        return s.length() > 3;
    });
// 此时不会输出

// 终端操作：触发执行
stream.collect(Collectors.toList());
// 此时才会输出过滤过程
				      \end{lstlisting}
				            \item \textbf{惰性求值机制}：
				                  \begin{itemize}[label={}, nosep]
					                  \item 中间操作构建操作链
					                  \item 终端操作触发整个链的执行
					                  \item 只在需要时才执行（优化性能）
				                  \end{itemize}
				            \item \textbf{代码验证}：
				                  \begin{lstlisting}[language=Java]
List<String> list = Arrays.asList("A", "BB", "CCC", "DDDD");

// 创建流（惰性）
Stream<String> stream = list.stream()
    .filter(s -> s.length() > 2)
    .map(String::toUpperCase);

// 触发执行
List<String> result = stream.collect(Collectors.toList());
System.out.println(result); // [CCC, DDDD]
				      \end{lstlisting}
			            \end{itemize}

			      \item[B] \textbf{forEach()是终端操作 — 正确}
			            \begin{itemize}[label={}, leftmargin=*, nosep]
				            \item \textbf{API定义}：
				                  \begin{lstlisting}[language=Java]
// Stream接口中定义
void forEach(Consumer<? super T> action);

// 返回void，说明是终端操作
				      \end{lstlisting}
				            \item \textbf{代码验证}：
				                  \begin{lstlisting}[language=Java]
List<String> list = Arrays.asList("A", "B", "C");

// forEach是终端操作
list.stream()
    .map(String::toUpperCase)
    .forEach(System.out::println);

// 可以添加peek（中间操作）在forEach之前
list.stream()
    .peek(System.out::println)  // 中间操作
    .map(String::toUpperCase)
    .forEach(System.out::println); // 终端操作
				      \end{lstlisting}
				            \item \textbf{常见终端操作}：
				                  \begin{itemize}[label={}, nosep]
					                  \item forEach
					                  \item collect
					                  \item count
					                  \item reduce
					                  \item anyMatch/allMatch/noneMatch
					                  \item findFirst/findAny
					                  \item min/max
				                  \end{itemize}
			            \end{itemize}

			      \item[C] \textbf{终端操作只能执行一次 — 正确}
			            \begin{itemize}[label={}, leftmargin=*, nosep]
				            \item \textbf{错误示例}：
				                  \begin{lstlisting}[language=Java]
Stream<String> stream = list.stream()
    .filter(s -> s.length() > 2);

// 第一次消费
long count1 = stream.count();

// 第二次消费：抛出IllegalStateException
// long count2 = stream.count(); // 错误！流已被消费
				      \end{lstlisting}
				            \item \textbf{异常信息}：
				                  \begin{lstlisting}[language=Java]
Exception in thread "main" java.lang.IllegalStateException:
stream has already been operated upon or closed
				      \end{lstlisting}
				            \item \textbf{设计原因}：
				                  \begin{itemize}[label={}, nosep]
					                  \item Stream是单次使用的管道
					                        \终端操作后Stream被消费
					                        \避免重复计算，提高性能
				                  \end{itemize}
				            \item \textbf{正确做法}：
				                  \begin{lstlisting}[language=Java]
// 方案1：每次创建新流
list.stream().filter(...).collect(...);
list.stream().filter(...).count();

// 方案2：将结果存储到变量
List<String> filtered = list.stream()
    .filter(s -> s.length() > 2)
    .collect(Collectors.toList());
// 可以多次使用结果
filtered.size();
filtered.forEach(...);
				      \end{lstlisting}
			            \end{itemize}

			      \item[D] \textbf{终端操作后可以继续进行中间操作 — 错误（本题答案）}
			            \begin{itemize}[label={}, leftmargin=*, nosep]
				            \item \textbf{错误示例}：
				                  \begin{lstlisting}[language=Java]
Stream<String> stream = list.stream()
    .map(String::toLowerCase);

// 执行终端操作
stream.forEach(System.out::println);

// 尝试添加中间操作：编译错误
// stream.filter(s -> s.length() > 2); // 错误！
                      \end{lstlisting}
				            \item \textbf{正确的操作顺序}：
				                  \begin{lstlisting}[language=Java]
// 正确：中间操作 终端操作
list.stream()
    .filter(s -> s.length() > 2)  // 中间
    .map(String::toUpperCase)     // 中间
    .forEach(System.out::println); // 终端

// 错误：终端操作 中间操作
Stream<String> s = list.stream().filter(...);
s.forEach(...);
// s.filter(...); // 错误
				      \end{lstlisting}
				            \item \textbf{原理机制}：
				                  \begin{itemize}[label={}, nosep]
					                  \item Stream构建：中间操作链接
					                  \item Stream执行：终端操作触发
					                  \item Stream关闭：终端操作后被消耗
					                  \item 无法在关闭的Stream上进行操作
				                  \end{itemize}
				            \item \textbf{命题陷阱}：这是Stream最重要的生命周期特性
			            \end{itemize}
		      \end{itemize}

		      \textbf{【应背诵知识点}
		      \begin{itemize}[label={}, nosep]
			      \item 中间操作：惰性、可链式调用、返回Stream
			      \item 终端操作：触发执行、只能执行一次、返回具体结果
			      \item Stream生命周期：构建→执行→消耗
			      \item 重要原则：终端操作后Stream不可再使用
		      \end{itemize}

		      \textbf{【命题趋势与实战策略】}

		      \textbf{考场秒杀}：看到"终端操作后可以继续..."→必错，Stream只能使用一次
	      \end{bluebox}

\end{enumerate}

\end{document}